\section{Descriptive Statistics}

The descriptive section is organized around the empirical facts that motivate the
regression analysis. We first define the working sample, then document where
workers are located in the firm-size distribution, how wages vary across that
distribution, and how much of the raw pattern is related to formality and other
observable wage gaps.

%\begin{table}[htbp]
  \centering
  \caption{Sample construction and coverage}
  \label{tab:sample-construction}
  \small
  \begin{tabular}{lp{0.58\textwidth}}
    \toprule
    Statistic & Value \\
    \midrule
    Harmonized worker-year file & 3,921,364 \\
    Valid real hourly income and positive weights & 3,921,364 \\
    Regression/descriptive sample & 3,852,322 \\
    Years covered & 2008--2025 (17 survey years; 2020 absent) \\
    Firm-size categories & Solo, 2-3, 4-5, 6-10, 11-19, 20-30, 31-50, 51-100, 101+ \\
    \bottomrule
  \end{tabular}
\end{table}


%Table~\ref{tab:sample-construction} defines the working sample and period coverage. The sample contains just over 5.0 million worker-year observations with valid wages, firm-size information, worker controls, and expansion weights over 17 survey years. Detailed employment, wage, and worker-composition tables are reported in Appendix~\ref{app:additional-descriptives}; the main text keeps the focus on the exhibits that structure the interpretation.

Figure~\ref{fig:income-mean-comparison} establishes the first raw fact behind the paper: mean real hourly income rises sharply with firm size. The gradient is visible in both endpoint years. In 2008, solo workers earned 5,431 COL 2025 pesos per hour on average, compared with 13,026 COL 2025 pesos per hour among workers in firms with 101 or more employees, a ratio of 2.4 to 1. By 2025, the corresponding means were 6,534 and 16,859 COL 2025 pesos per hour, a ratio of 2.6 to 1. %The middle categories also sit well above solo work in 2025: workers in firms with 6--10 employees earned 9,663 COL 2025 pesos per hour, and those in firms with 51--100 employees earned 13,217 COL 2025 pesos per hour. 
The bubble sizes show that these are not comparisons among negligible cells: solo work accounted for 34.7 percent of employment in 2025, while the 101+ category accounted for 25.2 percent.

\begin{figure}[htbp]
  \centering
  \includegraphics[width=0.95\textwidth]{fig65_\figurelang.png}
  \caption{Mean real hourly income by firm size}
  \label{fig:income-mean-comparison}
\end{figure}

Figure~\ref{fig:employment-share-comparison} establishes the second raw fact: employment is heavily concentrated at the bottom of the firm-size distribution. In 2025, solo workers accounted for 34.7 percent of employment, firms with 2--3 workers for 13.8 percent, firms with 4--5 workers for 7.0 percent, and firms with 6--10 workers for 5.8 percent. Taken together, solo work and firms with up to ten workers represented 61.2 percent of employment. By contrast, the middle of the distribution was thin: firms with 11--100 workers jointly accounted for 13.6 percent. The largest category grew substantially, from 18.8 percent of employment in 2008 to 25.2 percent in 2025, but this increase did not eliminate the large employment mass in solo work and micro-firms, which still covered almost three out of every five workers in 2025.

\begin{figure}[htbp]
  \centering
  \includegraphics[width=0.9\textwidth]{fig64_\figurelang.png}
  \caption{Employment distribution by firm size: 2008 and 2025}
  \label{fig:employment-share-comparison}
\end{figure}

Figure~\ref{fig:wage-gaps-2025} puts the firm-size gradient in context by comparing it with other raw remuneration gaps in 2025. The gender gap is small in these weighted means: men earn 0.6 percent below the overall mean and women 0.8 percent above it. The formal-informal gap is much larger. Formal workers earn 15,241 COL 2025 pesos per hour, 48.0 percent above the overall mean, while
informal workers earn 6,271 COL 2025 pesos per hour, 39.1 percent below it; the formal-informal wage ratio is therefore about 2.4 to 1, close to the 2.6-to-1 ratio between 101+ firms and solo work in
Figure~\ref{fig:income-mean-comparison}. Education and sector gaps are also large: workers with higher education earn 71.6 percent above the overall mean, while workers with preschool or less earn 57.9 percent below it; across activities, agriculture is 41.4 percent below the mean, whereas information and communication is 112.8 percent above it. The raw firm-size gap is therefore far
larger than the gender gap, comparable to the formality gap, and embedded in education and sectoral sorting. Appendix Table~\ref{tab:descriptive-wage-gaps-2025} reports the corresponding means, employment shares, and gaps.

\begin{figure}[!htbp]
  \centering
  \includegraphics[width=0.95\textwidth]{fig73_\figurelang.png}
  \caption{Mean real hourly income gaps by sex, formality, education, and sector, 2025; gaps are relative to the overall 2025 mean. The sector panel omits extraterritorial organizations.}
  \label{fig:wage-gaps-2025}
\end{figure}

Appendix Table~\ref{tab:descriptive-wage-groups-mean-2008-2025} reports mean real hourly
wages in 2008 and 2025 for all workers, by sex, and by formality status. The
table gives the exact magnitudes behind Figure~\ref{fig:income-mean-comparison}.
Three patterns stand out. First, the raw firm-size gradient was already large in
2008 and remained large in 2025: mean hourly income in the 101+ category was
about 2.4 times the solo-worker mean in 2008 and about 2.6 times the
solo-worker mean in 2025. Second, the gender gap in the aggregate means is small
relative to the cross-firm-size gradient. Third, the formal-informal wage gap is
large in both years, making formality a central part of the interpretation.
Appendix Table~\ref{tab:descriptive-wage-growth-2008-2025} reports the
corresponding wage changes between 2008 and 2025.

\ifpolicyreport
\begin{table}[htbp]
  \centering
  \caption{Real hourly labor income by firm size, 2025}
  \label{tab:descriptive-income-size-2025}
  \small
  \begin{tabular}{lrrrrrr}
    \toprule
    Firm size & Mean & P25 & Median & P75 & S.D. & Raw premium \\
    \midrule
    Solo & 8,351 & 3,606 & 5,769 & 8,942 & 11,807 & 0.0\% \\
    2-3 & 8,830 & 4,615 & 6,522 & 8,654 & 13,608 & 5.7\% \\
    4-5 & 9,503 & 5,440 & 6,923 & 8,654 & 16,211 & 13.8\% \\
    6-10 & 10,609 & 6,393 & 7,466 & 9,532 & 17,310 & 27.0\% \\
    11-19 & 12,346 & 6,844 & 7,867 & 11,276 & 18,438 & 47.8\% \\
    20-30 & 11,997 & 7,139 & 8,145 & 11,538 & 13,414 & 43.6\% \\
    31-50 & 13,428 & 7,141 & 8,515 & 13,516 & 17,573 & 60.8\% \\
    51-100 & 14,481 & 7,212 & 8,897 & 14,685 & 18,799 & 73.4\% \\
    101+ & 18,502 & 7,615 & 11,037 & 20,769 & 21,986 & 121.5\% \\
    \bottomrule
  \end{tabular}
  \vspace{0.3em}
  \begin{minipage}{0.95\textwidth}
  \footnotesize
  Notes: Incomes use only 2025 observations and are expressed as real hourly labor income in constant 2025 pesos. All statistics are weighted by GEIH expansion weights. The raw premium is the percent difference in the weighted mean relative to solo workers in 2025.
  \end{minipage}
\end{table}


\begin{table}[htbp]
  \centering
  \caption{Worker characteristics by firm size, 2025}
  \label{tab:descriptive-worker-characteristics-2025}
  \small
  \begin{tabular}{lrrrr}
    \toprule
    Firm size & Women & Formal & Higher education & Mean log wage \\
    \midrule
    Solo & 46.2\% & 14.1\% & 25.2\% & 8.639 \\
    2-3 & 42.4\% & 20.1\% & 26.1\% & 8.769 \\
    4-5 & 42.6\% & 32.2\% & 28.3\% & 8.879 \\
    6-10 & 44.0\% & 54.4\% & 36.3\% & 9.009 \\
    11-19 & 44.2\% & 70.4\% & 46.1\% & 9.128 \\
    20-30 & 45.1\% & 84.2\% & 49.5\% & 9.168 \\
    31-50 & 44.2\% & 91.9\% & 56.3\% & 9.254 \\
    51-100 & 46.8\% & 94.8\% & 58.4\% & 9.311 \\
    101+ & 47.5\% & 98.3\% & 67.6\% & 9.511 \\
    \bottomrule
  \end{tabular}
  \vspace{0.3em}
  \begin{minipage}{0.95\textwidth}
  \footnotesize
  Notes: Statistics use only 2025 observations and are weighted by GEIH expansion weights. Higher education corresponds to workers classified as superior or university educated.
  \end{minipage}
\end{table}


For the policy report, Tables~\ref{tab:descriptive-income-size-2025} and
\ref{tab:descriptive-worker-characteristics-2025} summarize the latest observed
labor-market structure in 2025.
\fi

\ifpolicyreport
\else
Figure~\ref{fig:income-formality-comparison} separates the endpoint comparison
by formality status. The bubble size represents the percentage of formal or
informal workers in each firm-size category in the corresponding year. Formal
workers earn more than informal workers and are much more concentrated in large
firms, but the wage gradient by firm size is visible within both formality
groups. This is why formality is a central mechanism but not the only object of
interest.

\begin{figure}[!htbp]
  \centering
  \includegraphics[width=0.95\textwidth]{fig69_\figurelang.png}
  \caption{Mean real hourly income by firm size and formality status, 2008 and 2025; dashed hollow markers denote 2008, solid filled markers denote 2025, and bubble size shows each firm-size category's employment share within the corresponding formality-year}
  \label{fig:income-formality-comparison}
\end{figure}

Figure~\ref{fig:income-formality-comparison} also shows that the firm-size gradient is not simply a formal-informal contrast. Among informal workers, mean hourly income still rises from 5,677 pesos for solo workers to 12,813 pesos in firms with 101 or more workers in 2025. Among formal workers, however, the profile is closer to a U-shape: formal solo workers earned 15,990 pesos per hour in 2025, formal workers in firms with 4--5 workers earned 12,036 pesos, and formal workers in firms with 101 or more workers earned 16,944 pesos. A plausible interpretation is that formal solo work is a selected category that includes independent professionals, such as lawyers, physicians, and consultants. This interpretation is consistent with the high earnings of formal solo workers and with their relatively small employment share within formal employment.

Figures~\ref{fig:income-sex-comparison} and
\ref{fig:income-education-comparison} extend the endpoint comparison to sex and
education groups.

\begin{figure}[!htbp]
  \centering
  \includegraphics[width=0.95\textwidth]{fig70_\figurelang.png}
  \caption{Mean real hourly income by firm size and sex, 2008 and 2025; dashed hollow markers denote 2008, solid filled markers denote 2025, and bubble size shows each firm-size category's employment share within the corresponding sex-year}
  \label{fig:income-sex-comparison}
\end{figure}

Figure~\ref{fig:income-sex-comparison} shows that the raw firm-size gradient is
similar for men and women. In 2025, male solo workers earned 6,816 pesos per
hour, compared with 16,929 pesos in firms with 101 or more workers; for women,
the corresponding means were 6,133 and 16,777 pesos. The large-firm advantage is
therefore visible within both sexes. The figure also shows that the gender gap
within a given firm-size category is small relative to the gap between solo work
and large-firm employment.

\begin{figure}[!htbp]
  \centering
  \includegraphics[width=\textwidth]{fig72_\figurelang.png}
  \caption{Mean real hourly income by firm size and education group, 2008 and 2025; dashed hollow markers denote 2008, solid filled markers denote 2025, and bubble size shows each firm-size category's employment share within the corresponding education-year}
  \label{fig:income-education-comparison}
\end{figure}

Figure~\ref{fig:income-education-comparison} is especially informative because
it separates two mechanisms that are otherwise mixed together. First, education
composition matters. In 2025, workers with preschool or less, primary, and
secondary education were much more concentrated in solo work than workers with
higher education. Solo work accounted for 65.2 percent of employment among
workers with preschool or less and 55.5 percent among workers with primary
education, compared with only 17.3 percent among workers with higher education.
Conversely, the 101+ category accounted for 47.3 percent of workers with higher
education but only 1.4 percent of workers with preschool or less and 4.9 percent
of workers with primary education. This sorting helps explain why the raw
firm-size wage gradient is so steep.

Second, education composition does not exhaust the pattern. Within every
education group, workers in the largest firms earned more than solo workers in
2025. The ratio between mean hourly income in 101+ firms and solo work was 2.1
among workers with preschool or less, 1.8 among workers with primary education,
1.5 among workers with secondary education, and 1.7 among workers with higher
education. At the same time, the higher-education panel has a distinctive shape:
solo workers with higher education earned more than workers in firms with 2--3
employees, before the gradient rose toward large firms. This is consistent with
the interpretation that some solo work among highly educated workers reflects
professional self-employment rather than only low-productivity informality. The
analogous comparison by sector is reported in
Appendix~\ref{app:additional-descriptives}.
\fi

\ifpolicyreport
\noindent Figures~\ref{fig:policy-employment-income-2025} and
\ref{fig:policy-composition-2025} summarize the 2025 policy snapshot, combining
employment, wage, formality, and education gradients by firm size.

\begin{figure}[htbp]
  \centering
  \includegraphics[width=0.9\textwidth]{fig66_\figurelang.png}
  \caption{Employment and mean hourly income by firm size in 2025}
  \label{fig:policy-employment-income-2025}
\end{figure}

\begin{figure}[htbp]
  \centering
  \includegraphics[width=0.9\textwidth]{fig67_\figurelang.png}
  \caption{Formality and higher education by firm size in 2025}
  \label{fig:policy-composition-2025}
\end{figure}
\fi

Taken together, the descriptive evidence motivates the regression analysis in
four ways. First, employment is concentrated in low-scale units: in the pooled
regression sample, solo workers and firms with up to ten workers account for
65.3 percent of weighted employment; in 2025 they still account for 61.3
percent. Second, the middle of the firm-size distribution is thin: firms with
11--100 workers account for only 12.3 percent of pooled employment. Third, the
largest harmonized category is economically important: firms with 101 or more
workers account for 22.5 percent of pooled employment. Fourth, the raw wage
gradient is steep, but larger firms also employ more formal and more educated
workers. The econometric exercise below therefore asks how much of the raw
gradient survives after controlling for observable worker composition,
formality, and sector-year conditions.

%\clearpage
